\chapter{Obiettivi e domande di ricerca}
\label{cap:obiettivi}

\section{Obiettivo}

Valutare, con esperimenti ripetibili, quanto le scelte architetturali proteggono un servizio dai guasti, e confrontare i risultati con previsioni formulate prima delle misure.

\section{Domande di ricerca}

\begin{table}[htbp]
\centering
\caption{Le domande di ricerca e le versioni confrontate.}
\label{tab:rq}
\begin{tabularx}{\linewidth}{@{}lLll@{}}
\toprule
\textbf{ID} & \textbf{Domanda} & \textbf{Tema} & \textbf{Confronto} \\
\midrule
RQ1 & Quanto fermo si evita passando da una a più copie dell'API? & ridondanza senza stato & V1 / V2 \\
RQ2 & Cosa cambia, in tempi, ordini persi e ripristino, inserendo una coda tra API e database? & asincronia & V2 / V3 \\
RQ3 & Quanto dipende il ripristino dal tipo di guasto e dai parametri delle sonde? & rilevamento & tutte \\
RQ4 & Quanto fermo evita una replica del database con failover, e quanti dati si rischiano? & ridondanza con stato & V2 / V4 \\
RQ5 & L'autoscaling sulla lunghezza della coda dà un compromesso migliore tra coda massima e costo, durante un picco? & elasticità & V3 / V5 \\
RQ6 & Quali risultati locali si ripetono su AWS, e perché gli altri cambiano? & trasferibilità & locale / AWS \\
\bottomrule
\end{tabularx}
\end{table}

\section{Ipotesi di partenza}

Le ipotesi qualitative sono elencate qui; le previsioni numeriche sono nel capitolo~\ref{cap:previsioni}.

\begin{enumerate}[label=\textbf{H\arabic*}]
  \item Più copie dell'API riducono molto il fermo dovuto all'API, ma non quello dovuto al database, che resta unico (RQ1).
  \item Con la coda il servizio continua ad accettare ordini anche con il database fermo e li salva al suo ritorno; in cambio cresce il tempo tra ricezione e salvataggio (RQ2).
  \item Un processo bloccato è rilevato più tardi di uno terminato, e il ritardo dipende dalla frequenza delle sonde (RQ3).
  \item Con l'arresto del solo processo, cioè il guasto F4 che misuriamo, la replica del database non accorcia necessariamente il fermo: nodo e disco restano integri, quindi il riavvio sul posto è già rapido quanto un failover. Il vantaggio atteso è nella perdita dell'intero nodo del database: senza replica PostgreSQL non può ripartire, perché i dati sono sul disco di quel nodo (D0006), mentre lo standby su un altro nodo può subentrare. Questo secondo caso non si misura ed è un limite dichiarato. Con replica asincrona si possono perdere gli ultimi ordini, con replica sincrona ogni scrittura è più lenta (RQ4).
  \item L'autoscaling smaltisce un picco con meno risorse di una configurazione fissa sovradimensionata, ma con un ritardo iniziale (RQ5).
\end{enumerate}

\section{Fuori ambito}

Non fanno parte del progetto: la replica del broker, la replica del control plane di Kubernetes, un database distribuito su più zone del cloud, la sicurezza e l'ottimizzazione delle prestazioni del codice. Il control plane è unico e non viene mai guastato (capitolo~\ref{cap:validita}).
